\section{箱线图与全部原始样本}
\begin{example}[text and listing,sidebyside=false]{同源样本，不分箱、不抽样}
\begin{gradtable}[style=rug-report,width=\linewidth]
\gradcolumn[width=25mm]{batch}{Batch}
\gradcolumn[type=boxrug,scale=fixed,
 options={min=0,max=20,height=6mm,label=none}]
 {samples}{Box plot and raw rug}
\gradrow{Atlas,{0,2,2,4,8,16,NA}}
\gradrow{Birch,{0,0,0,NA}}
\gradrow{Cedar,{2,3,4,5,6}}
\gradrow{Dune,{NA,NA}}
\end{gradtable}
\end{example}
独立调用 \key{\gradboxrug[options]{raw samples}}；默认高 5ex。每条短线对应一个有效观测，重复位置可重叠但重数不变。沿用 type-7 分位数、1.5-IQR 实际观测须和离群点，不抖动、不抽样、不自动分箱。最多接收 2048 个原始位置（含 NA），不接受预计算摘要。
零、单样本、常量、全缺失沿用既有统计；固定、自动、命名范围与裁剪不改变统计。一个组合只生成一个 Figure，包含统计和保留重复值的完整有效样本列表。表格样式 \key{rug-report} 和图形样式 \key{sample-rug} 均为显式选择。
